% Generated draft: requires mathematical and source review.
\documentclass[11pt]{article}
\usepackage[margin=25mm]{geometry}
\usepackage{fontspec}
\setmainfont{Latin Modern Roman}
\usepackage{amsmath,amssymb,mathtools,unicode-math}
\setmathfont{Latin Modern Math}
\usepackage{enumitem,needspace,xurl,hyperref}
\setlength{\parindent}{0pt}
\setlength{\parskip}{5pt}
\setlength{\emergencystretch}{3em}
\newcommand{\N}{\mathbb{N}}
\newcommand{\Z}{\mathbb{Z}}
\newcommand{\Q}{\mathbb{Q}}
\newcommand{\R}{\mathbb{R}}
\newcommand{\C}{\mathbb{C}}
\newcommand{\F}{\mathbb{F}}
\newcommand{\A}{\mathbb{A}}
\newcommand{\PP}{\mathbb{P}}
\newcommand{\E}{\mathbb{E}}
\newcommand{\eps}{\varepsilon}
\newcommand{\dd}{\,\mathrm{d}}
\begin{document}
\section*{UnsolvedMath ID: 800717}
\textbf{Source problem code:} catalogue entry
\subsection{Definitions and mathematical statement}
Let \(F\) be a field, and let \(A\) and \(B\) be finite-dimensional central division algebras over \(F\), using the central simple algebra context of the supplied source. For a central simple \(F\)-algebra \(C\), write \(C\cong M_r(D)\), where \(D\) is a central division \(F\)-algebra, and define its index by \(\operatorname{ind}(C)=\sqrt{\dim_F D}\). In particular, put \(a=\operatorname{ind}(A)\) and \(b=\operatorname{ind}(B)\).

A field extension \(E/F\) is a common splitting field of \(A\) and \(B\) if
\[
A\otimes_F E\cong M_a(E)
\qquad\text{and}\qquad
B\otimes_F E\cong M_b(E)
\]
as \(E\)-algebras.

\textbf{Problem 31.} Suppose that
\[
\operatorname{ind}(A\otimes_F B)
<\operatorname{ind}(A)\operatorname{ind}(B).
\]
Describe the common splitting fields of \(A\) and \(B\).

\textbf{Unresolved scope.} The supplied formulation asks what these common splitting fields are, but specifies neither the desired form of a characterization nor any degree bound. It does not say whether the intended extensions must be finite, separable, or Galois. The definition above applies to any field extension; no additional restriction on the requested splitting fields is inferred. The source introduces this question as an alternative to considering maximal subfields, but the preceding formulation is not supplied.
\subsection{Short English statement}
For two central division algebras over the same field, suppose the index of their tensor product is strictly smaller than the product of their indices. The problem asks for a description of fields that split both algebras. The intended restrictions on those fields and the precise kind of description sought remain unspecified.
\subsection{Sources}
\begin{itemize}
\item Local UnsolvedMath registry, permanent ID \(800717\), exact source problem code \texttt{catalogue entry}, titled \emph{Open Problems on Central Simple Algebras --- Problem 31}. The supplied record attributes the problem to Auel, Brussel, Garibaldi, and Vishne, with proposed year 2010.
\item Auel, Brussel, Garibaldi, and Vishne, \emph{Open Problems on Central Simple Algebras} (2010), Problem 31: \url{https://arxiv.org/abs/1006.3304}. The supplied provenance records access on 2026-10-16 and extraction by \texttt{arxiv\_latex}; these are supplied provenance statements, not independently verified access or source checks.
\item The supplied registry labels this problem \texttt{open} and records an update on 2026-10-16. This is supplied dated curation, not an independently verified assessment of the mathematical literature or current status.
\end{itemize}
\end{document}
